% ===========================================================================
% ASPLOS 2027 -- separate note describing changes since previous submissions.
% Uploaded to the submission site alongside the paper.  One page, per the CFP.
% This document is reviewer-visible, so it is anonymous like the paper.
% ===========================================================================
\documentclass[10pt]{article}
\usepackage[letterpaper,margin=0.9in]{geometry}
\usepackage{enumitem}
\usepackage{microtype}
\setlist[itemize]{leftmargin=*,itemsep=1.0pt,topsep=2pt,parsep=0pt}
\pagestyle{empty}

\begin{document}

\begin{center}
{\large\textbf{Changes Since Previous Submissions}}\\[0.3em]
\emph{Compression Ratio Is Not a Serving Metric: A Measurement Study of
KV-Cache Compression for Long-Context LLM Serving}
\end{center}
\vspace{0.3em}

\noindent\textbf{Previous submissions.}
This work was submitted twice to \emph{Transactions on Machine Learning
Research} (TMLR) during 2026. The first submission was titled after the
evaluation harness itself; we do not reproduce that title here, because the
harness is public under that name and quoting it would deanonymise this
submission. The second was titled ``Beyond Compression Ratio: Evaluating the
Algorithmic and Systems Utility of Long-Context KV-Cache Compression''.
The submission identifiers are available to the chairs on request.
Both were desk-rejected by the editors, and in neither case was the paper sent
out for review: \textbf{no external reviewer reports were produced and no
technical feedback of any kind was received.} The notices stated that the
submission was judged unlikely to meet one or both of TMLR's criteria and could
not be advanced given submission volume and limited reviewer and action-editor
capacity.

\smallskip
\noindent We state this plainly because it determines how this document should
be read: \textbf{none of the changes below responds to a review}, because there
was none. They are our own reassessment of how the work should be framed and
for which community. Both earlier versions were written as descriptions of an
evaluation framework; we now think the framework was the least interesting
thing about the results, and that the work belongs at a systems venue as a
measurement study.

\medskip
\noindent\textbf{What changed.}
\begin{itemize}
\item \textbf{Re-centred from a benchmark to a systems finding}, with a new
  title, abstract and contribution ordering. The paper leads with a measured
  claim about serving --- that compression ratio and task score are not
  sufficient to predict prefill latency, decode throughput and resident memory
  --- and introduces the methodology only after that claim is established;
  both earlier versions introduced the framework first. The contributions now
  lead with the empirical characterisation rather than the composite scores.
\item \textbf{The first two pages were rebuilt around that finding}, in
  recognition of the rapid-review round: by the end of page~2 the reader has
  the systems problem, why conventional reporting is insufficient for it, the
  three quantitative findings, and their consequences for capacity planning
  and latency objectives; the explicit contribution bullets follow immediately
  on page~3.
\item \textbf{A conceptual problem in the metric exposition was fixed.} The
  earlier versions argued that representation utility and serving efficiency
  are distinct axes while reporting a composite $\mathcal{S}_{\rm sys}$ that
  contains $\mathcal{S}_{\rm res}$ by construction. The pure hardware term
  $\phi$ is now the serving axis, $\mathcal{S}_{\rm sys}$ is stated to be a
  presentational convenience and never evidence, and the reported rank
  displacements are identified as a lower bound. No formula and no measured
  value changed.
\item \textbf{Statistical language was scoped to the evidence}: claims of the
  form ``none predicts the others'' became statements about what is and is not
  detectable at this sample size, while direct matched observations are stated
  as the direct observations they are.
\item \textbf{Related work was rewritten for a systems audience}, and now
  engages concurrent work that reaches a compatible headline conclusion from a
  different measurement design, stating what our design adds to it.
\item \textbf{Length was reduced to 11 pages by content editing}, not by
  spacing or font changes: redundant exposition was cut and secondary material
  (proofs, alternative formulations, per-task and per-context detail, protocol,
  artifact appendix) moved to appendices.
\item \textbf{Figures were regenerated at final printed size} so every label
  renders at 8\,pt or larger without scaling, and captions now state the
  systems takeaway rather than describe the axes.
\item \textbf{An artifact appendix was added}, and the manuscript was
  anonymised for double-blind review: the harness is referred to by a pseudonym
  because it is public under its own name, and prior work by the authors is
  cited in the third person.
\end{itemize}

\medskip
\noindent\textbf{What did not change.}
\textbf{No new experiments were run and no measured value was altered.} The
study is the same 3{,}150 paired query executions over the same seven methods,
three models, five tasks and context lengths to 128K. Every number, table and
figure is regenerated from the same released per-query records, and a
verification script re-derives each quantity quoted in the prose and fails if
the manuscript has drifted from the data. The limitations --- single GPU, batch
size one, 2B--7B models, and the implementation caveats for the two methods
whose production kernels were unavailable on our platform --- are unchanged and
remain stated in the body.

\end{document}
